\section{总控 Skill：Using Superpowers}

\subsection{总控流程}

\texttt{Using Superpowers} 是第一个被触发的 Skill（"use when starting any conversation"）。它定义了整体开发链路：

\begin{enumerate}
    \item \textbf{Brainstorming}（头脑风暴）：与用户多轮对话确认修改边界。
    \item \textbf{Writing Plans}（写实施计划）：将设计转化为可执行步骤。
    \item \textbf{Development}（开发）：以 TDD 方式先写测试再写实现。
    \item \textbf{Verification}（验证）：完成前真实运行命令检查结果。
\end{enumerate}

\begin{importantbox}{Skill.md 中的强制触发指令}
Superpowers 的 \texttt{skill.md} 中包含用 XML 标签包裹的强制指令：\textit{"极度重要——如果你认为有 1\% 的机会可以使用某个 Skill，一定要触发它。"}
\end{importantbox}

\subsection{流程的形式化定义}

Skill 内部使用 \textbf{Dot 语言}（Graphviz）定义流程图——用形式化代码语言描述 SOP，使语言模型更易消化理解。流程中还要求 Agent \textbf{宣称}（declare）自己正在使用什么 Skill、出于什么目的，并定义一组 \textbf{Checklist} 检查是否完成要求。

\subsection{本章小结}

Using Superpowers 是 Agent 的"操作系统入口"，通过 Dot 语言定义的有向图来约束 Agent 行为，并通过宣称机制和 Checklist 确保每一步都有迹可查。


